% Draft results section template — placeholders to be filled once the
% Phase-A/B-refined 5-drone A/B/C/D campaign and the MC inter-fault-gap
% sweep complete.
%
% A post-hoc script `fill_results_placeholders.py` (TODO) can sed the
% \NEWRmsA ... macros with actual numerical values once data is ready.

\section{Results}
\label{sec:results}

% =====================================================================
% Placeholder macros for numbers pending the Phase-A/B campaign re-run.
% =====================================================================
\newcommand{\NEWRmsA}{\ensuremath{\langle\!\langle\mathrm{RMS}_A\rangle\!\rangle}}
\newcommand{\NEWPeakA}{\ensuremath{\langle\!\langle\mathrm{Peak}_A\rangle\!\rangle}}
\newcommand{\NEWPeakTA}{\ensuremath{\langle\!\langle T^{\max}_A\rangle\!\rangle}}
\newcommand{\NEWPeakFA}{\ensuremath{\langle\!\langle F^{\max}_A\rangle\!\rangle}}
\newcommand{\NEWSigA}{\ensuremath{\langle\!\langle\sigma_{T,A}\rangle\!\rangle}}
\newcommand{\NEWRmsB}{\ensuremath{\langle\!\langle\mathrm{RMS}_B\rangle\!\rangle}}
\newcommand{\NEWPeakB}{\ensuremath{\langle\!\langle\mathrm{Peak}_B\rangle\!\rangle}}
\newcommand{\NEWPeakTB}{\ensuremath{\langle\!\langle T^{\max}_B\rangle\!\rangle}}
\newcommand{\NEWPeakFB}{\ensuremath{\langle\!\langle F^{\max}_B\rangle\!\rangle}}
\newcommand{\NEWSigB}{\ensuremath{\langle\!\langle\sigma_{T,B}\rangle\!\rangle}}
\newcommand{\NEWRmsC}{\ensuremath{\langle\!\langle\mathrm{RMS}_C\rangle\!\rangle}}
\newcommand{\NEWPeakC}{\ensuremath{\langle\!\langle\mathrm{Peak}_C\rangle\!\rangle}}
\newcommand{\NEWPeakTC}{\ensuremath{\langle\!\langle T^{\max}_C\rangle\!\rangle}}
\newcommand{\NEWPeakFC}{\ensuremath{\langle\!\langle F^{\max}_C\rangle\!\rangle}}
\newcommand{\NEWSigC}{\ensuremath{\langle\!\langle\sigma_{T,C}\rangle\!\rangle}}
\newcommand{\NEWRmsD}{\ensuremath{\langle\!\langle\mathrm{RMS}_D\rangle\!\rangle}}
\newcommand{\NEWPeakD}{\ensuremath{\langle\!\langle\mathrm{Peak}_D\rangle\!\rangle}}
\newcommand{\NEWPeakTD}{\ensuremath{\langle\!\langle T^{\max}_D\rangle\!\rangle}}
\newcommand{\NEWPeakFD}{\ensuremath{\langle\!\langle F^{\max}_D\rangle\!\rangle}}
\newcommand{\NEWSigD}{\ensuremath{\langle\!\langle\sigma_{T,D}\rangle\!\rangle}}
\newcommand{\NEWGapRatio}{\ensuremath{\langle\!\langle T_C/T_D\rangle\!\rangle}}
\newcommand{\NEWQPactive}{\ensuremath{\langle\!\langle n_{\mathrm{act}}\rangle\!\rangle}}
\newcommand{\NEWQPsolve}{\ensuremath{\langle\!\langle t_{\mathrm{qp}}\rangle\!\rangle}}

\subsection{Experimental Setup}
\label{sec:results:setup}
The controller, plant, fault model, and integration scheme are those
specified in Section~\ref{sec:methods}. All experiments use five
quadrotors ($m_i=\SI{1.5}{\kilo\gram}$) lifting a rigid payload
($m_L=\SI{3.0}{\kilo\gram}$) through $L=\SI{1.25}{\metre}$ cables
discretised into $N_b=8$ beads. The reference trajectory is a 3-D
Bernoulli lemniscate of amplitude $\SI{3}{\metre}$, vertical
oscillation $A_z=\SI{0.35}{\metre}$, and period $T=\SI{12}{\second}$.
Each scenario runs for $\SI{40}{\second}$ at
$\Delta t=\SI{0.2}{\milli\second}$. Scenario~A is fault-free; B severs
drone~0 at $t=\SI{15}{\second}$; C adds a drone~2 severance at
$t=\SI{20}{\second}$; D delays the second severance to
$t=\SI{25}{\second}$.

\subsection{Nominal Tracking}
\label{sec:results:nominal}
Scenario~A establishes the fault-free baseline. Following the Phase-A/B
initialisation fix (Section~\ref{sec:hover-ic}), the chaotic pickup
transient present in earlier campaigns is eliminated, yielding
$\mathrm{RMS}\,\lVert e_p\rVert=\NEWRmsA~\si{\metre}$ with a peak of
$\NEWPeakA~\si{\metre}$ (Table~\ref{tab:results}). Peak cable tension
is $\NEWPeakTA~\si{\newton}$ and per-drone thrust peaks at
$\NEWPeakFA~\si{\newton}$. Fig.~F02 plots the three-dimensional
payload track against the reference; Fig.~F03 overlays the
tracking-error time series; Fig.~F05 shows the load-sharing
imbalance $\sigma_T(t)\approx 0$ throughout cruise, evidencing
perfectly symmetric steady-state operation.

\subsection{Single-Fault Response}
\label{sec:results:single}
In Scenario~B, drone~0's cable is severed at $t=\SI{15}{\second}$. The
remaining four agents redistribute the load through the
$T_{\mathrm{ff}}=T_{\mathrm{meas}}$ feed-forward identity, without any
centralised reallocation. Post-fault
$\mathrm{RMS}\,\lVert e_p\rVert=\NEWRmsB~\si{\metre}$, peak
$\NEWPeakB~\si{\metre}$, peak tension $\NEWPeakTB~\si{\newton}$.
Fig.~F04 (per-rope tension waterfall) shows the severed rope's
segments collapsing to zero tension in one control tick while the
surviving ropes' segments absorb the displaced load;
Fig.~F06 (STFT of $\sigma_T$) shows the fault event as a broadband
vertical stripe, uniquely detectable from one scalar signal.

\subsection{Compound-Fault Timing Effect}
\label{sec:results:compound}
Scenarios~C and~D isolate the effect of inter-fault gap. With a
$\SI{5}{\second}$ gap (C), the second severance lands while the first
transient is still settling: peak tension reaches
$\NEWPeakTC~\si{\newton}$ --- approximately $\NEWGapRatio$ higher
than the $\SI{10}{\second}$-gap case D ($\NEWPeakTD~\si{\newton}$),
consistent with constructive overlap of the two rebalancing
transients. Notably, tracking error is essentially invariant across
B--D ($\NEWRmsB$, $\NEWRmsC$, $\NEWRmsD~\si{\metre}$), demonstrating
that emergent rebalancing is robust to inter-fault timing even when
the peak-tension envelope is not.

\begin{table}[t]
\centering
\caption{5-Drone campaign metrics (Phase-A/B refreshed).}
\label{tab:results}
\begin{tabular}{lccccc}
\toprule
Scenario & RMS $\lVert e_p\rVert$ & Peak $\lVert e_p\rVert$ & Peak $T$ & Peak $\lVert F\rVert$ & RMS $\sigma_T$ \\
 & [m] & [m] & [N] & [N] & [N] \\
\midrule
A nominal     & \NEWRmsA & \NEWPeakA & \NEWPeakTA & \NEWPeakFA & \NEWSigA \\
B single      & \NEWRmsB & \NEWPeakB & \NEWPeakTB & \NEWPeakFB & \NEWSigB \\
C dual 5\,s   & \NEWRmsC & \NEWPeakC & \NEWPeakTC & \NEWPeakFC & \NEWSigC \\
D dual 10\,s  & \NEWRmsD & \NEWPeakD & \NEWPeakTD & \NEWPeakFD & \NEWSigD \\
\bottomrule
\end{tabular}
\end{table}

\subsection{Monte-Carlo Inter-Fault-Gap Sweep}
\label{sec:results:sweep}
A parametric sweep over inter-fault gap
$\Delta t\in\{2,4,6,8,10,12,14,16,18,20\}~\si{\second}$ isolates the
worst-case transient envelope. Fig.~F10 scatters peak rope tension
against $\Delta t$ with a running-max envelope, confirming the
monotone bound from Scenario~D to the adversarial-worst case near
$\Delta t = \SI{2}{\second}$, beyond which the two transients
decouple in time.

\subsection{Instrumentation Diagnostics}
\label{sec:results:diag}
QP active-set occupancy (Fig.~F07) is zero across all scenarios during
cruise; the box-constraints activate only during the final descent
segment, supporting the claim that the controller operates in an
unconstrained convex regime during normal operation. Median QP solve
time is $\NEWQPsolve~\si{\micro\second}$, well below the $200~\si{\micro\second}$
budget of the $\Delta t = 0.2~\si{\milli\second}$ control period.
Fig.~F09 places the closed-loop payload-vertical poles in the left
half plane for every $N_{\mathrm{alive}}\in\{1,\dots,5\}$, certifying
internal stability down to the single-surviving-drone limit.

The thrust and tilt envelope timelines (Fig.~F08) show
$\lVert F_i\rVert$ peaking at $\NEWPeakFA~\si{\newton}$ during the
lift phase, well below $F_{\max}=\SI{150}{\newton}$, and per-segment
rope tension waterfalls during pickup (Fig.~S01) confirm the absence
of the snap-taut impulsive transient that characterised the previous
linear-ramp initial condition.
